\documentclass[11pt,a4paper]{article}

% This file is deliberately compatible with the default pdfLaTeX compiler on
% Overleaf/ShareLaTeX.  The accompanying zip contains the figures in figures/.
\usepackage[a4paper,margin=25mm]{geometry}
\usepackage[T1]{fontenc}
\usepackage[utf8]{inputenc}
\usepackage{lmodern}
\usepackage{microtype}
\usepackage{amsmath,amssymb}
\usepackage{booktabs}
\usepackage{tabularx}
\usepackage{graphicx}
\usepackage{caption}
\usepackage{subcaption}
\usepackage{float}
\usepackage{placeins}
\usepackage{enumitem}
\usepackage{xcolor}
\usepackage{hyperref}

\definecolor{navy}{HTML}{17365D}
\definecolor{blue}{HTML}{2F6690}
\hypersetup{
  colorlinks=true,
  linkcolor=navy,
  urlcolor=blue,
  pdftitle={Neural Operators for 21cmFAST Lightcone Regeneration},
  pdfauthor={Haris}
}

\setlength{\parindent}{0pt}
\setlength{\parskip}{0.55em}
\setlist[itemize]{leftmargin=1.5em,itemsep=0.2em,topsep=0.2em}
\setlist[enumerate]{leftmargin=1.7em,itemsep=0.2em,topsep=0.2em}
\captionsetup{font=small,labelfont=bf}
\renewcommand{\arraystretch}{1.15}

\newcommand{\xhi}{x_{\mathrm{HI}}}
\newcommand{\Ltwo}{L^2}
\newcommand{\Hone}{H^1}
\newcommand{\FNO}{\textsc{FNO}}
\newcommand{\UFNO}{\textsc{U-FNO}}
\newcommand{\WHNO}{\textsc{WHNO}}
\newcommand{\code}[1]{\texttt{#1}}

\title{\textbf{Neural Operators for 21cmFAST Lightcone Regeneration}\\
\large From Fourier baselines to generalized bases and local--global operators}
\author{Haris}
\date{August 2026}

\begin{document}
\maketitle
\thispagestyle{empty}

\begin{abstract}
This report develops and evaluates neural-operator surrogates for regenerating
21cmFAST neutral-fraction lightcones from matter-density fields and the
simulation parameters that determine their reionization histories. The work
follows a clear architectural progression. A conditioned Fourier Neural
Operator (\FNO) establishes a useful global baseline; a U-shaped Fourier
operator (\UFNO) then adds a local real-space path; and a generalized operator
matrix tests Fourier, SIREN-generated Fourier, Walsh--Hadamard, wavelet, and
convolutional choices independently in local and global slots.

Parameter conditioning reduces validation \(\Ltwo\) from about \(0.20\) to
\(0.06\), showing that the density field alone does not identify the timing of
reionization. The \UFNO subsequently reaches validation
\((\Ltwo,\Hone)=(0.0418,8.27)\), improving on the pure-\FNO result. In the
final 3-D matrix, \code{sfno/swhno} comes within \(2.3\%\) of \UFNO RMSE with
271 times fewer parameters, while \code{cnn/whno} is the fastest accurate
model. The Walsh--Hadamard transform transfers from the 2-D experiments to
the 3-D lightcone task, particularly in the global slot. The results point to
a practical design principle: global operators determine large-scale
communication, while the local path has the greater influence on sharp
structure and computational cost.
\end{abstract}

\tableofcontents
\clearpage

\section{Introduction}

The goal is to replace the expensive forward calculation of a 21cmFAST
lightcone with a learned field-to-field map,
\begin{equation}
  \mathcal{G}:\bigl(\delta_m(\mathbf{x},z),\theta\bigr)
  \longmapsto \widehat{\xhi}(\mathbf{x},z),
  \label{eq:map}
\end{equation}
where \(\delta_m\) is the matter-density lightcone and \(\theta\) contains the
cosmological and astrophysical parameters used to generate the target. The
output is a three-dimensional neutral-fraction field rather than a single
summary statistic. This makes the problem useful as a test of neural operators
in a setting with both long-range correlations and thin, spatially organized
ionization fronts. Fourier neural operators provide the global spectral
construction used here \cite{li2021fno}.

The experiments were not designed as a search over arbitrary network
architectures. They follow a sequence of increasingly specific questions:
\begin{enumerate}
  \item Can a global Fourier operator learn the lightcone map at all?
  \item Which input information is needed to determine the reionization
  history?
  \item Does changing the representation of the global operator help?
  \item Does a local path improve the reconstruction more than additional
  global spectral capacity?
\end{enumerate}
The report concentrates on the positive results that answer these questions.
Negative comparisons are included only where they explain why the next
architectural step was taken.

\section{Data, task, and evaluation}

\subsection{Lightcones and inputs}

The production data contain \(6600\) 21cmFAST lightcones
\cite{mesinger2010}. The cached cubes have
shape \(140\times140\times256\) and cover \(z=5\)--\(25\). The input to the
production models has thirteen channels:
\begin{itemize}
  \item the matter-density field, scaled by a factor of ten;
  \item \(1/(1+z)\), broadcast along the volume; and
  \item eleven cosmological and astrophysical parameters, standardized using
  training-set statistics and broadcast over the volume.
\end{itemize}
The target is \(\xhi\) on the same grid. Splits are made at the cone level,
so chunks from one lightcone cannot occur in both training and validation
sets. This matters because neighboring chunks share both density structure and
reionization history.

The main runs use distributed data parallel training on four H200 GPUs. The
per-rank batch size is one. Adam with cosine annealing is used throughout the
baseline comparisons. The original cache is uniform in redshift; this report
does not include the later line-of-sight grid experiments, since they change
the sampled target rather than the operator architecture.

\subsection{Metrics}

The primary optimization metrics are the absolute field error \(\Ltwo\) and
the Sobolev error \(\Hone\). The held-out analysis adds diagnostics that
measure different aspects of the regenerated field:
\begin{itemize}
  \item per-cone RMSE and \(R^2\), separating overall timing from spatial
  structure;
  \item the global neutral-fraction history \(\overline{\xhi}(z)\);
  \item the power-spectrum ratio and the cross-correlation \(r(k)\);
  \item the mean-free-path bubble-size distribution; and
  \item the width and error of the transition band around an ionization front.
\end{itemize}
No single metric describes this task. A model can obtain a good volume-averaged
error while still smoothing the fronts, because most of the volume is close to
the saturated values zero or one.

\section{A conditioned Fourier baseline}

\subsection{Why the parameters are part of the input}

The first useful baseline is a standard spectral operator applied to the
density field. A density-only model converges to a nearly constant prediction:
different lightcones can have similar density structure but different
reionization histories. This is an identifiability problem, not simply a
training failure.

Broadcasting the eleven simulation parameters changes the problem. The model
can now distinguish cones with similar density fields but different global
histories. Validation \(\Ltwo\) falls from approximately \(0.20\) for the
density-only input to approximately \(0.06\) with conditioning, while
\(\Hone\) falls from \(17.85\) to about \(11.5\). The first major gain is
therefore supplied by information in the input rather than by a larger
operator.

\subsection{Fourier layers}

An \FNO layer combines a pointwise map with a truncated Fourier convolution,
\begin{equation}
  v^{(\ell+1)}(\mathbf{x}) =
  \sigma\!\left(
  W^{(\ell)}v^{(\ell)}(\mathbf{x})+
  \mathcal{F}^{-1}\!\left[
    R^{(\ell)}(\mathbf{k})\,
    \mathcal{F}[v^{(\ell)}](\mathbf{k})
  \right](\mathbf{x})
  \right).
  \label{eq:fno-layer}
\end{equation}
Here \(R^{(\ell)}\) is learned only for the retained modes. The spectral path
communicates across the full field in one layer, which makes it a natural
starting point for a lightcone with large bubbles and long-range timing
correlations.

The conditioned pure \FNO reaches a test \(\Ltwo\) of about \(0.057\) and a
test \(\Hone\) of \(11.64\) in the main spectral comparison. Increasing the
retained modes from \(16^3\) to \(24^3\) does not improve the converged
solution. This result motivates changing the representation and the
architecture, rather than continuing to increase the Fourier cutoff.

\begin{figure}[t]
\centering
\includegraphics[width=0.92\textwidth]{figures/fno_training_trajectories.png}
\caption{Training trajectories for the principal lightcone models. Conditioning
and the local path produce the clearest improvements over the initial global
Fourier baseline.}
\label{fig:trajectories}
\end{figure}

\section{Generalized operator bases}

\subsection{Measuring the learned spectral preference}

The nominal mode cutoff does not tell us which modes the trained network uses.
To measure this directly, the RMS weight of each spectral mode was logged
throughout training. In a 32-mode \UFNO, the lowest eight modes eventually
carry about \(52\%\) of the spectral weight, compared with \(25\%\) for a
uniform profile. The same qualitative preference appears in a 64-mode model.

\begin{figure}[t]
\centering
\includegraphics[width=0.84\textwidth]{figures/mode_weight_profiles_3d.png}
\caption{Mode-weight profiles from the 3-D operator diagnostics. The nominal
bandwidth is larger than the effective bandwidth used by the trained Fourier
kernel.}
\label{fig:mode-weights}
\end{figure}

This measurement distinguishes two ideas that are often conflated. A model
may have enough modes to represent a sharp feature while still assigning most
of its learned weight to low modes. In this task, the output of the nonlinear
network is not determined by the retained Fourier modes alone. The diagnostic
therefore provides a basis for comparing different operator representations
without treating the mode count as a complete explanation.

\subsection{SIREN-generated Fourier kernels}

SirenFNO \cite{shi2025siren} replaces the table of independent weights with a sinusoidal
hypernetwork that generates the Fourier kernel as a function of mode
coordinates. In the tested implementation, the generator is evaluated across
the full mode grid rather than stopping at a conventional low-frequency
cutoff. The number of generator parameters is independent of the resolution
at which the kernel is queried.

The learned spectrum remains approximately flat in the 3-D experiment. More
importantly, this representation improves the plain Fourier baseline:
\begin{table}[H]
\centering
\caption{Held-out comparison of the principal spectral representations.}
\label{tab:siren}
\begin{tabular}{@{}lrr@{}}
\toprule
Model & Test \(\Ltwo\) & Test \(\Hone\)\\
\midrule
Plain FNO & 0.0570 & 11.64\\
SirenFNO & 0.0501 & 9.82\\
\UFNO & \textbf{0.0397} & \textbf{7.87}\\
\bottomrule
\end{tabular}
\end{table}

SirenFNO therefore demonstrates that full-frequency kernel parameterization is
useful for this map. It does not, by itself, reach the \UFNO result. The
comparison is informative: the SIREN model changes the global spectral
representation, whereas the \UFNO also supplies a local real-space route.

\subsection{Walsh--Hadamard operators}

The Walsh--Hadamard transform (WHT) \cite{perez2025whno} provides a different
generalization of the spectral basis. Its basis functions are rectangular
\(\pm1\) step waves, and
the fast transform has \(O(n\log n)\) complexity using additions and
subtractions. Coefficients are ordered by sequency, the step-wave analogue of
frequency.

The WHT is especially attractive for a target containing extended regions
separated by sharp fronts. In the 2-D slice experiments, replacing the global
Fourier operator with \WHNO gives the best result when the local slot remains
Fourier-based. The improvement comes from the global slot; placing the WHT
only in the local window does not provide the same benefit.

The final 3-D matrix confirms that the global WHT transfers beyond the 2-D
task. At a matched local operator, \code{whno/swhno} reaches RMSE \(0.0654\),
compared with \(0.0741\) for \code{whno/whno}. The stronger Walsh-global
pairings outperform the all-Fourier \code{fno/fno} reference in the matrix.
This is the first indication in the project that a generalized basis can
improve the global part of the lightcone operator without increasing its
parameter count.

\section{The local--global operator design}

\subsection{The U-FNO result}

The \UFNO combines global Fourier blocks with a small 3-D U-Net branch,
following the U-shaped operator constructions described by Wen et al.
\cite{wen2022ufno} and Rahman et al. \cite{rahman2023uno}. The Fourier path
supplies long-range communication,
while strided convolutions,
upsampling, and skip connections provide a local route for high-resolution
features. With the same conditioned inputs and training objective as the pure
\FNO, the \UFNO reaches
\begin{equation}
\boxed{\mathrm{val}\ \Ltwo=0.0418,\qquad
\mathrm{val}\ \Hone=8.27.}
\end{equation}
Its test \(\Ltwo\) is \(0.0408\) in the corresponding benchmark. This is a
25--29\% improvement over the main pure-\FNO result and is the strongest
single architectural improvement in the baseline sequence.

\begin{figure}[p]
\centering
\includegraphics[width=0.78\textwidth]{figures/lightcone_grid_3d_validation.png}
\caption{Representative validation lightcones for the \UFNO. The model
reproduces global timing and bubble morphology; the visible residuals are
concentrated near transition regions.}
\label{fig:ufno-grid}
\end{figure}

\begin{figure}[p]
\centering
\begin{subfigure}[t]{0.48\textwidth}
\centering
\includegraphics[width=\linewidth]{figures/comparison_3d_validation_cone61.png}
\caption{Heavily reionized cone.}
\end{subfigure}\hfill
\begin{subfigure}[t]{0.48\textwidth}
\centering
\includegraphics[width=\linewidth]{figures/comparison_3d_validation_cone2322.png}
\caption{Patchy transition.}
\end{subfigure}
\caption{Two slices from held-out \UFNO predictions. Large-scale structure and
the location of the transition are reproduced across different histories.}
\label{fig:ufno-slices}
\end{figure}
\FloatBarrier

\subsection{Separating the local and global roles}

The final matrix treats the operator as two slots: a local operator acts in
overlapping windows and a global operator acts on the bottleneck. Fourier,
SIREN-Fourier, WHT, SIREN-WHT, wavelet, and convolutional operators can be
placed in either slot. This setup makes it possible to ask whether the basis
or the spatial role is responsible for an observed change.

The held-out leaderboard is shown in Table~\ref{tab:leaderboard}. The \UFNO
remains the most accurate model, but several much smaller local--global
operators are close to it.

\begin{table}[H]
\centering
\caption{Best held-out 3-D results for the main operator combinations. The
parameter counts treat a complex Fourier weight as two real parameters.}
\label{tab:leaderboard}
\small
\begin{tabular}{@{}lrrr@{}}
\toprule
Local / global & Parameters & RMSE & \(R^2\)\\
\midrule
\UFNO & 202,902,721 & \textbf{0.05781} & \textbf{0.9508}\\
\code{sfno/swhno} & 747,841 & 0.05914 & 0.9485\\
\code{cnn/whno} & 3,468,081 & 0.06134 & 0.9446\\
\code{cnn/swhno} & 1,405,921 & 0.06184 & 0.9437\\
\code{whno/swhno} & 1,050,241 & 0.06544 & 0.9369\\
\code{swhno/swhno} & 677,857 & 0.06632 & 0.9352\\
\code{whno/whno} & 2,609,137 & 0.07405 & 0.9192\\
\code{fno/whno} & 5,705,713 & 0.07594 & 0.9150\\
\code{fno/fno} & 20,385,777 & 0.07989 & 0.9060\\
\code{wno/whno} & 2,182,129 & 0.08933 & 0.8824\\
\bottomrule
\end{tabular}
\end{table}

Two results stand out. First, \code{sfno/swhno} is within \(2.3\%\) of the
\UFNO RMSE with 271 times fewer parameters. This is a substantial storage
reduction, although the SIREN kernel regeneration makes this model slower at
inference. Second, the convolutional local slot is consistently strong:
\code{cnn/whno} is close to \UFNO accuracy while using only \(1.7\%\) of its
parameters.

\begin{figure}[t]
\centering
\includegraphics[width=0.78\textwidth]{figures/matrix_params_vs_rmse.png}
\caption{Parameter count and held-out RMSE for the final matrix. Several
generalized-basis models are close to \UFNO accuracy at much smaller size.}
\label{fig:params-rmse}
\end{figure}

\subsection{Inference cost}

Parameter count is not a reliable proxy for inference speed in this design.
On an A100, batch size one, and complete \(140\times140\times256\) cubes,
\code{cnn/whno} takes \(115.9\) ms per cube, compared with \(274.1\) ms for
\UFNO. Its peak memory is \(4345\) MiB, compared with \(10271\) MiB for
\UFNO. The very small \code{sfno/swhno} model takes \(458.1\) ms because the
SIREN generator rebuilds the kernel on each forward pass.

The global basis has almost no effect on the measured runtime: \code{fno/fno}
and \code{fno/whno} both take about \(347.1\) ms. Most of the cost is in the
local path, particularly the window loop for the windowed operators. The
practical choice therefore depends on the intended use. \UFNO gives the
lowest RMSE, \code{cnn/whno} gives the best speed--accuracy trade-off, and
\code{sfno/swhno} gives the smallest checkpoint among the near-best models.

\begin{figure}[t]
\centering
\includegraphics[width=0.78\textwidth]{figures/matrix_inference_speed_vs_rmse.png}
\caption{Inference speed versus RMSE for the final matrix. The two useful
speed--accuracy choices are \UFNO and the convolutional-local
\code{cnn/whno} model.}
\label{fig:speed-rmse}
\end{figure}

\subsection{Morphology and phase coherence}

Whole-volume RMSE does not fully rank the models by morphology. The matrix
analysis finds that the mean small-scale cross-correlation
\(\langle r(k)\rangle_{k>1}\) predicts the bubble-size bias with Spearman
\(\rho=+0.88\), whereas the small-scale power amplitude has no independent
relationship with that bias. This is useful for the operator comparison:
adding small-scale power is not enough if the added structure is out of phase
with the target.

The convolutional local models perform well on this morphological diagnostic.
For example, \code{cnn/swhno} has bubble-size bias \(+0.006\), while its RMSE
is \(0.0618\). The result supports the architectural split suggested by the
matrix: the local branch should preserve spatially coherent boundaries, while
the global basis should communicate the large-scale field.

\begin{figure}[t]
\centering
\includegraphics[width=0.72\textwidth]{figures/matrix_coherence_vs_bsd_bias.png}
\caption{Small-scale cross-correlation against bubble-size bias. Phase
coherence is more informative for this morphology statistic than power
amplitude alone.}
\label{fig:coherence-bsd}
\end{figure}

\section{Discussion}

The results give a consistent picture of what the neural operator must do for
this problem.
\begin{enumerate}
  \item \textbf{Condition on the simulation parameters.} The density field
  does not determine the reionization history by itself. The first large
  improvement comes from supplying the missing global information.
  \item \textbf{Use a global operator for communication.} Fourier and
  Walsh--Hadamard operators both provide a compact route across the full
  volume. The WHT result shows that the basis can be changed without changing
  the local--global skeleton.
  \item \textbf{Add a local path for detailed structure.} The \UFNO
  improvement is larger than the gains from increasing the Fourier cutoff.
  The local convolutional path is the clearest architectural explanation for
  the reduction in field and gradient error.
  \item \textbf{Evaluate beyond voxel error.} Bubble morphology follows
  small-scale phase coherence more closely than power amplitude. A model
  intended for downstream physical analysis should therefore be checked in
  both field space and structure space.
\end{enumerate}

The strongest immediate model depends on the constraint. \UFNO is the
accuracy reference. \code{cnn/whno} is attractive when throughput and memory
matter. \code{sfno/swhno} shows that a generalized, generated basis can
approach the accuracy of a much larger model with a very small parameter
count, although its current implementation does not deliver the same
inference efficiency. These are complementary outcomes rather than a single
universal ranking.

\section{Conclusion}

The project has established a useful neural-operator path from a conditioned
global \FNO to a generalized local--global design. Conditioning makes the
lightcone map learnable. SIREN-generated Fourier kernels improve the pure
spectral baseline by distributing the learned representation across the
available modes. Walsh--Hadamard operators transfer successfully from 2-D to
3-D, with their clearest benefit in the global slot. Finally, the \UFNO
demonstrates that a local real-space route is more valuable than simply
increasing global Fourier capacity for the sharp, structured parts of the
field.

The final matrix makes the trade-offs concrete. The \UFNO provides the best
RMSE, \code{sfno/swhno} nearly matches it with 271 times fewer parameters, and
\code{cnn/whno} is the most attractive speed--accuracy compromise. Together,
these results support a simple design for future work: retain a global
generalized operator for long-range structure, pair it with an efficient local
path, and select models using morphology and throughput as well as voxelwise
error.

\begin{thebibliography}{9}
\bibitem{li2021fno}
Z. Li et al.,
\emph{Fourier Neural Operator for Parametric Partial Differential Equations},
International Conference on Learning Representations, 2021.

\bibitem{rahman2023uno}
M. A. Rahman, Z. E. Ross, and K. Azizzadenesheli,
\emph{U-NO: U-shaped Neural Operators},
Transactions on Machine Learning Research, 2023.

\bibitem{wen2022ufno}
G. Wen, Z. Li, K. Azizzadenesheli, A. Anandkumar, and S. M. Benson,
\emph{U-FNO---An enhanced Fourier neural operator-based deep-learning model
for multiphase flow},
Advances in Water Resources, 163, 104180, 2022.

\bibitem{shi2025siren}
P. Shi, J. Yin, S. Tierney, and J. Gao,
\emph{SirenFNO: Efficient and Full Frequency Learning of Fourier Neural Operators},
2025.

\bibitem{perez2025whno}
M. Pérez Cuadrado et al.,
\emph{Walsh-Hadamard Neural Operators for Solving PDEs with Discontinuous
Coefficients}, arXiv:2511.07347, 2025.

\bibitem{mesinger2010}
A. Mesinger, S. Furlanetto, and R. Cen,
\emph{21cmFAST: A Fast, Semi-Numerical Simulation of the 21 cm Signal from the
Epoch of Reionization}, Monthly Notices of the Royal Astronomical Society,
411, 955--972, 2011.
\end{thebibliography}

\end{document}
